\chapter{Theoretical Fundamentals}
\label{c:fundamentals}

%%%%%%%%%%%%%%%%%%%%%%%%%%%%%%%%%%%%%%%%%%%%%%%%%%%%%%%%%%%%%%%%%%%%%%%%%%%%%%%%

% ==========================================
% Section: Fermionic Many-Body Systems
% ==========================================
\section{Fermionic Many-Body Systems}
Fermions, such as electrons, protons, and neutrons, are particles and primary building blocks of matter.
They differ from their properties and interactions between each other.
A Many-Body System describes these interactions between the particles.
Specifically, a simulation must describe these rules and properties as accurately as possible.
\cite[chapter 6.1]{benito_kottmann_qc_2025}

% ==========================================
% Subsection: Indistinguishability
% ==========================================
\subsection{Indistinguishability}
In classical Physics, identical objects and particles can be easily tracked and measured in quantum mechanics this is not possible anymore.
Identical particles get completely indistinguishable, which means we cannot label the particles nor can we measure their location over time, because this would collapse or disturb the systems state.
Follows, that if two of the systems particles get swapped the system stays in the same state as before.
This is because physical quantities are calculated as expectation values ($\bra{\psi} H \ket{\psi} = \sum_{i}{\abs{c_i}^2} E_i$) and the probability amplitude stays the same.
\cite[chapter 7.1]{sakurai2020modern}\cite[chapter 6.1]{benito_kottmann_qc_2025}


% ==========================================
% Subsection: Permutation Symmetry
% ==========================================
\subsection{Permutation Symmetry}
A system containing only identical particles can either be totally symmetric or antisymmetric under the interchange of any pair.
For Fermions this interchange would always be strictly antisymmetric, which introduces a global phase factor of $-1$ to the wavefunction, when applying such an interchange.
Bosons on the other hand are strictly symmetric and therefore do not apply any phase at all.
\cite[chapter 7.2]{sakurai2020modern}~\cite[chapter 6.1]{benito_kottmann_qc_2025}

% ==========================================
% Subsection: Pauli Exclusion Principle
% ==========================================
\subsection{Pauli Exclusion Principle}\label{subsec:pauli-exclusion-principle}
We have now enough information to describe the Pauli Exclusion Principle.
It states that `Two Fermions can not be in the same state.'
We can see this easily when investigating such a system.
Let's take two Fermions (let us call them $\ket{\psi}$), then the quantum state for the system can be described as $\ket{\psi} \otimes \ket{\psi}$.
This state can not be antisymmetric so it can not include Fermions.
\cite[chapter 7.1, 7.2]{sakurai2020modern}~\cite[chapter 6.1]{benito_kottmann_qc_2025}

% ==========================================
% Section: The Second Quantization Framework
% ==========================================
\section{The Second Quantization Framework}\label{sec:the-second-quantization-framework}
In second Quantization we introduce operators to deal with the antisymmetrical property of Fermions.
We no longer need to handle these properties in a complex way, but rather bake them into an annihilation ($\hat{a}^\dagger_k$) and creation ($\hat{a}_k$) operator.
These two operators annihilate/create an electron at spin-orbital $k$.
Furthermore, to preserve the antisymmetric wavefunctions these operators anticommute.
This implies some physical boundaries:
\begin{align}
    \hat{a} |0\rangle &= 0, \\
    \hat{a}^\dagger |0\rangle &= |1\rangle, \\
    \hat{a} |1\rangle &= |0\rangle, \\
    \hat{a}^\dagger |1\rangle &= 0.
\end{align}
From the Pauli Exclusion Principle also follows that we cannot annihilate two electrons with the same spin at the same position (because only one can exist there), and also we can not create more than one electron with the same spin at the same position. \\
To understand these very important operators more we can hav a look at the Jordan Wigner (JW) Transformation, that maps fermionic spin orbitals to qubits.
This transformation uses non-local Pauli strings to construct composite operators:
\begin{align}
    \label{eq:creation_op_jd_transformation}
    \hat{a}_k^\dagger &= \mathbb{I}^{\otimes (k-1)} \otimes \hat{\sigma}_k^+ \otimes \hat{\sigma}_z^{\otimes (N-k)}, \\
    \label{eq:annhilation_op_jd_transformation}
    \hat{a}_k &= \mathbb{I}^{\otimes (k-1)} \otimes \hat{\sigma}_k^- \otimes \hat{\sigma}_z^{\otimes (N-k)}.
\end{align}
All orbitals before the $k$-th where must be left untouched which here is defined by $\mathbb{I}^{\otimes (k-1)}$.
The annihilation and creation takes place in the local ladder operators, which are defined as,
\begin{equation}
    \hat{\sigma}_k^\pm = \frac{1}{2}(\hat{\sigma}_x \pm i\hat{\sigma}_y).\label{eq:ladder_operators}
\end{equation}
The non-local parity string $\hat{\sigma}_z^{\otimes (N-k)}$ applies a phase of $-1$ for every occupied orbital and therefore enforces the Pauli exclusion principle.

We now defined new operators to deal with the more complex properties of Fermions, we still need a way to store these states instead of complex antisymmetrical wavefunctions, without being restricted to a fixed numer of particles.
A standard one particle system is described in a Hilbert space, we can use a sum of these Hilbert spaces to build a new space.
We call it the Fock Space.\\
Within the Fock Space we can express a state without the explicit spatial coordinates of every particle, instead we are using the occupation number representation:
\begin{equation}
    \ket{\psi} = \ket{n_1, n_2, \dots, n_i, \dots}.\label{eq:occupation_number_representation}
\end{equation}
In this notation, each Index $i$ represents a single-particle state (or spin-orbital), while $n_i$ denotes its respective occupation number.
Due to the Pauli Exclusion Principle the occupation numbers for Fermions are strictly restricted to $n_i \in \{0, 1\}$.

\cite[chapter 2.1]{schleich2025quantum}\cite[chapter 6.5]{benito_kottmann_qc_2025}\cite[chapter 7.7.1]{sakurai2020modern}
% JW transformation
\cite[chapter III.B]{anand2022quantum}\cite[chapter II.E]{kottmann2020feasible}\cite[chapter 1.1, 5.3]{Kottmann2025}

% ==========================================
% Subsection: Fermionic Excitations
% ==========================================
\subsection{Fermionic Excitations}
To dynamical evolve a quantum state, we need a tool to modify the positions of the electrons.
A so-called excitation does exactly that, it can change the position of one or more electrons within the state.
In quantum chemistry, a parameterized unitary excitation can be represented as,
\begin{equation}
    U(\theta) = e^{-i \frac{\theta}{2} G},
\end{equation}
where $G$ is the Generator.
It stores the information about the physical movement of the changing electrons. \\
We can now have a look at different Generators, to understand how these excitations work and what they achieve.
A single excitation can be expressed in the $G_{pq}$-Generator:
\begin{equation}
    G_{pq} = i(\hat{a}^\dagger_p\hat{a}_q - \hat{a}^\dagger_q\hat{a}_p).
\end{equation}
The Generator is Hermitian, due to the hermitian conjugate, the minus sign in the parentheses and the multiplication by $i$.
It ensures that the transition is reversible and probability preserving.
This Generator describes the movement of an electron from spin-orbital $p$ to $q$. \\

For more complex excitations we can use a $n$-folded Generator.
It describes $n$ single excitation chained together in one single Generator, which equals $n$ simultaneous transitions from occupied spin-orbitals $\mathbf{p} = \{p_1, \dots, p_n\}$ to targets $\mathbf{q} = \{q_1, \dots, q_n\}$.
We call it the $G_{\boldsymbol{pq}}$-Generator, it is defined as
\begin{equation}
    \label{eq:mult_fe_exc_generator}
    G_{\mathbf{pq}} = i \left( \prod_{k=1}^{n} \hat{a}^\dagger_{p_k} \hat{a}_{q_k} - \text{h.c.} \right),
\end{equation}
where h.c. denotes the Hermitian conjugate.
Also, the before discussed single excitation follows as a special case where $n = 1$.
\cite[Ch.~II]{kottmann2020feasible}\cite[Ch.~III.A]{anand2022quantum}.


















